\clearpage\thispagestyle{empty}
\pdfbookmark[0]{Quellen zum Bericht über Vesta vom Juli 1807}{vesta1807july-sources}
\begin{center}\textbf{Quellen zum Bericht über Vesta vom Juli 1807}\end{center}
\noindent Diese KI-Übertragung in heutiges Deutsch folgt Franz Xaver von Zachs Bericht vom Juli 1807 in \href{https://gdz.sub.uni-goettingen.de/id/PPN236007467}{Gauß’ \textit{Werken}, Band~VI} (Göttingen, 1874), S.~288--290. Der Bericht enthält Carl Friedrich Gauß’ Brief vom 7. Juni 1807 und seinen Vergleich der Beobachtungen mit den zweiten Bahnelementen der Vesta. Der Erstdruck erschien als Beitrag VII, \textit{Fortgesetzte Nachrichten über den neuen von Dr. Olbers in Bremen entdeckten Planeten Vesta}, in der \textit{Monatlichen Correspondenz zur Beförderung der Erd- und Himmelskunde}, Band~XVI, \href{https://www.digitale-sammlungen.de/en/view/bsb10538607?page=91}{S.~83}--\href{https://www.digitale-sammlungen.de/en/view/bsb10538607?page=100}{92}. Gauß’ Brief beginnt auf \href{https://www.digitale-sammlungen.de/en/view/bsb10538607?page=92}{S.~84}. Seine Vergleichstafel reicht von \href{https://www.digitale-sammlungen.de/en/view/bsb10538607?page=93}{S.~85} über \href{https://www.digitale-sammlungen.de/en/view/bsb10538607?page=94}{S.~86} bis \href{https://www.digitale-sammlungen.de/en/view/bsb10538607?page=95}{S.~87}. Dort folgen seine Bemerkungen zu drei Deklinationen.
\par\medskip
Die Tabelle der Werkausgabe enthält eine Spalte für die mittlere Zeit am Beobachtungsort. Die beiden Überschriften für die Unterschiede behalten die historischen Kürzel R.$-$B. und R$-$B. Der Auszug der Werkausgabe endet mit dem eingeklammerten Hinweis auf die Meridianbeobachtungen von Groombridge und Maskeline. Der Erstdruck gibt die entsprechenden Beobachtungen in Greenwich auf \href{https://www.digitale-sammlungen.de/en/view/bsb10538607?page=97}{S.~89} und schreibt den zweiten Namen dort Maskelyne. Die Zeitschriftendigitalisate stammen von der Bayerischen Staatsbibliothek. Übertragung und Satz wurden mit KI erstellt.
\clearpage
